\chapter{Admissible Sets and Compatibility Certificates}
\index{admissible set|textbf}
\index{compatibility certificate|textbf}

\label{ch:admissible-sets-certificates}

\chapterstatus{foundational model-adequacy framework with synthetic worked
examples; no material acceptance decision, physical threshold calibration, or
uncertainty-quantification result is claimed.}

A reduced Hamiltonian is rarely judged by one number. A candidate may reproduce
selected energy bands while failing to reproduce the represented operator, its
real-space locality, or an observable needed by the intended application. The
opposite can also occur: a candidate may be close in one operator norm while
missing a small but important spectral feature. This chapter develops a precise
way to ask whether one candidate satisfies two declared criteria at the same
time.

The central objects are \emph{admissible sets}. Each criterion selects the
parameters that satisfy its threshold. A common acceptable candidate exists
when the corresponding sets intersect. An explicit point in the intersection
is a \emph{common witness}. If the sets do not intersect, a valid positive
lower bound on their distance is a \emph{separation certificate}. A failed
numerical search is neither one of these.

After this chapter, the reader should be able to
\begin{itemize}
  \item distinguish model parameters from alignment variables and numerical
        controls;
  \item define normalized spectral and operator losses on declared data;
  \item distinguish established spectral fitting and operator regression from
        the proposed simultaneous witness-or-certificate question;
  \item construct thresholded admissible sets;
  \item distinguish a witness, a failed search, and a separation certificate;
  \item recognize ellipsoidal sublevel sets of positive-definite quadratic
        losses;
  \item keep training, withheld evaluation, and locality diagnostics in their
        proper roles; and
  \item state exactly what a bounded synthetic compatibility result does and
        does not establish.
\end{itemize}

\section{From one error to simultaneous requirements}
\index{model adequacy!simultaneous requirements}

Suppose that a reduced model is controlled by a parameter
$\boldsymbol\theta$ in a declared domain $\Theta$. A spectral comparison
produces a nonnegative loss $L_E(\boldsymbol\theta)$, while a represented-
operator comparison produces a nonnegative loss $L_H(\boldsymbol\theta)$. Two
thresholds, $\tau_E$ and $\tau_H$, state the largest losses accepted for the
present purpose.

It is not enough to minimize the losses separately. The minimizer of $L_E$ may
be a different model from the minimizer of $L_H$. The scientific question is
whether there is one parameter value for which
\begin{equation}
  L_E(\boldsymbol\theta)\leq\tau_E
  \quad\text{and}\quad
  L_H(\boldsymbol\theta)\leq\tau_H.
\end{equation}
The word ``and'' is the essential point. Reporting one spectrally successful
model and a different operator-successful model does not exhibit a single
adequate model.

\subsection{Example 1: intervals on a line}

Before considering Hamiltonians, take one real parameter $x$ and define
\begin{equation}
  L_E(x)=|x|,
  \qquad
  L_H(x)=|x-1|.
\end{equation}
With equal thresholds $\tau_E=\tau_H=0.6$, the two admissible sets are
\begin{equation}
  \mathcal A_E=[-0.6,0.6],
  \qquad
  \mathcal A_H=[0.4,1.6].
\end{equation}
Their intersection is $[0.4,0.6]$. The point $x=0.5$ is a common witness.
Neither $x=0$, which minimizes $L_E$, nor $x=1$, which minimizes $L_H$, is
needed to prove compatibility.

If both thresholds are instead $0.4$, then
\begin{equation}
  \mathcal A_E=[-0.4,0.4],
  \qquad
  \mathcal A_H=[0.6,1.4].
\end{equation}
The sets are separated by $0.2$. This positive number is stronger evidence
than checking a finite list of points and failing to find an intersection.

\section{The candidate family and its contracts}
\index{candidate family}
\index{parameter domain}
\index{model-class expressivity}

Let
\begin{equation}
  \mathfrak M
  =
  \left\{
    H_{\boldsymbol\theta}(k):
    \boldsymbol\theta\in\Theta
  \right\}
\end{equation}
be a family of finite Hermitian Bloch matrices. A complete definition of this
family states
\begin{itemize}
  \item the meaning, units, and allowed range of every component of
        $\boldsymbol\theta$;
  \item the reciprocal coordinates $k$ on which the matrices are evaluated;
  \item the retained rank, basis, gauge, energy reference, and energy unit;
  \item the parent or target data to which the candidates are compared; and
  \item the distinction between fitting data, withheld data, and purely
        diagnostic data.
\end{itemize}
A bounded domain is not merely a programming convenience. It is part of the
model class. Enlarging $\Theta$ asks a different question and may change both
compatibility and separation.

The matrix structure is likewise part of the scientific question, not merely
an implementation choice. Stiles showed how the independent matrix elements
allowed by a declared crystal structure can be enumerated before a tight-
binding fit~\cite{stiles1997}. Such enumeration determines which coordinates
are available to the fit; it does not prove that the resulting family can
reproduce every desired property. In a specified spin--orbit
$sp^3d^5s^{\ast}$ empirical tight-binding class, for example, Boykin, Klimeck,
and Oyafuso derived exact relations among selected band gaps and effective
masses~\cite{boykinKlimeckOyafuso2004}. Those relations are intrinsic
restrictions of that model class. This is stronger than observing that one
optimization run failed to locate a satisfactory parameter set.

\subsection{A rank-two shift-and-splitting family}
\index{rank-two model}

A useful controlled family begins with a rank-two reference matrix
$H_{\mathrm{ref}}(k)$. Define its mean energy by
\begin{equation}
  \overline E(k)
  =
  \frac{1}{2}\operatorname{tr}H_{\mathrm{ref}}(k)
\end{equation}
and its traceless part by
\begin{equation}
  H_{0}(k)
  =
  H_{\mathrm{ref}}(k)-\overline E(k)I_2.
\end{equation}
Given a positive energy scale $E_0$, consider
\begin{equation}
  H_{s,q}(k)
  =
  \overline E(k)I_2
  +sE_0I_2
  +qH_0(k).
  \label{eq:admissible-rank-two-family}
\end{equation}
The dimensionless parameter $s$ shifts both levels equally. The dimensionless
parameter $q$ rescales their splitting about the mean. The parameter column is
\begin{equation}
  \boldsymbol\theta
  =
  \begin{bmatrix}s\\q\end{bmatrix},
  \qquad
  \boldsymbol\theta\in\Theta.
\end{equation}
Equation~\eqref{eq:admissible-rank-two-family} is a deliberately restricted
model family. It does not represent every rank-two Hamiltonian. A compatibility
or separation result obtained from it is therefore conditional on this family
and its bounds.

When the domain requires $q>0$, the ordered eigenvalues obey
\begin{equation}
  E_n^{\mathrm{cand}}(k;s,q)
  =
  \overline E(k)+sE_0
  +q\left(E_n^{\mathrm{ref}}(k)-\overline E(k)\right).
  \label{eq:admissible-rank-two-eigenvalues}
\end{equation}
Each spectral residual is therefore affine in $(s,q)$. For any fixed global
alignment, the entries of the aligned candidate matrix are also affine in
$(s,q)$. Squaring and summing these residuals produces the quadratic losses
derived below. The positivity restriction on $q$ is important: allowing its
sign to change can reverse the ordering of the two eigenvalues and invalidate
this single ordered formula.

\section{Spectral and operator losses}
\index{loss!spectral}
\index{loss!operator}

The two criteria use the same declared sampling contract but compare different
objects. Let $k_1,\ldots,k_{N_k}$ be a frozen training mesh, let the retained
rank be $r$, and assign each sample a finite positive dimensionless weight
$w_j$. Write
\begin{equation}
  W=\sum_{j=1}^{N_k}w_j.
\end{equation}

\subsection{Spectral loss}
\index{spectral RMS error}
\index{normalized spectral loss}

Suppose the target energies are $E_n^{\mathrm{tar}}(k_j)$, ordered according
to a declared convention.

\begin{definition}[Spectral RMS error]
The energy-valued ordered spectral root-mean-square error of the candidate
specified by $\boldsymbol\theta$ is
\begin{equation}
  R_E(\boldsymbol\theta)
  =
  \left[
    \frac{1}{rW}
    \sum_{j=1}^{N_k}w_j\sum_{n=1}^{r}
    \left|
      E_n^{\mathrm{cand}}(k_j;\boldsymbol\theta)
      -E_n^{\mathrm{tar}}(k_j)
    \right|^2
  \right]^{1/2}.
  \label{eq:admissible-spectral-rms-error}
\end{equation}
It has the same energy unit as the candidate and target eigenvalues.
\end{definition}

\begin{definition}[Normalized spectral loss]
Given a finite positive energy scale $E_0$, the normalized spectral loss is
\begin{equation}
  L_E(\boldsymbol\theta)
  =
  \frac{R_E(\boldsymbol\theta)}{E_0}.
  \label{eq:admissible-spectral-rms}
\end{equation}
It is dimensionless.
\end{definition}

Uniform weights recover the factor $1/(N_k r)$. Another normalization defines
another loss and must be named rather than silently substituted. Neither the
weights nor $E_0$ are automatically uncertainty estimates.

Objectives built from energy residuals have a long history. Starrost and
collaborators used a genetic algorithm to fit an empirical tight-binding
Hamiltonian to energies at selected high-symmetry wave
vectors~\cite{starrost1996}. More recently, Wang and collaborators trained adjustable
tight-binding matrix elements with an eigenvalue mean-squared-error objective
against first-principles band energies~\cite{wangEtAl2021}. These are
precedents for spectral supervision, although their sampling, weights, and
normalizations need not equal those in
Equation~\eqref{eq:admissible-spectral-rms}.

The ordered-band formula also has a matching premise. Away from crossings, an
energy ordering may be sufficient. At crossings or inside nearly degenerate
clusters, the comparison must state how bands or clusters are identified. A
sorting rule, overlap-based tracking rule, projector comparison, or another
declared convention changes the loss contract and cannot be inferred from the
energy arrays alone.

\subsection{Operator loss}
\index{represented-operator RMS residual}
\index{normalized operator loss}

The corresponding operator comparison requires both matrices to act in an
identified retained space. Let $H^{\mathrm{tar}}(k_j)$ be the target matrices
in that space, and let $C$ be an allowed global unitary alignment.

\begin{definition}[Represented-operator RMS residual]
The energy-valued represented-operator root-mean-square residual of the
candidate specified by $\boldsymbol\theta$ under alignment $C$ is
\begin{equation}
  R_H(\boldsymbol\theta,C)
  =
  \left[
    \frac{1}{rW}
    \sum_{j=1}^{N_k}w_j
    \left\|
      C^{\dagger}H_{\boldsymbol\theta}(k_j)C
      -H^{\mathrm{tar}}(k_j)
    \right\|_{\mathrm F}^{2}
  \right]^{1/2}.
  \label{eq:admissible-operator-rms-residual}
\end{equation}
It has the same energy unit as the candidate and target matrices.
\end{definition}

\begin{definition}[Normalized operator loss]
Given the same finite positive energy scale $E_0$, the normalized operator
loss under alignment $C$ is
\begin{equation}
  L_H(\boldsymbol\theta,C)
  =
  \frac{R_H(\boldsymbol\theta,C)}{E_0}.
  \label{eq:admissible-operator-rms}
\end{equation}
It is dimensionless.
\end{definition}

The Frobenius norm and its sample aggregation are meaningful only after the
state spaces, bases, gauges, energy references, units, reciprocal coordinates,
sample identities, and weights have been declared. The scientific loss
definition chooses the samples and weights; represented-operator mechanics
validate compatibility, form differences, and apply the declared aggregation.
Chapter~\ref{ch:mathematical-foundations} develops those prerequisites.

Direct matrix supervision addresses a different object from spectral
supervision. Zhang and collaborators learned Hamiltonian and overlap blocks in
a declared atomic-orbital representation using equivariant atomic-cluster-
expansion models~\cite{zhangEtAl2022}. Schwade and collaborators instead
combined an approximate tight-binding model with learned environment-dependent
corrections and evaluated the resulting electronic-property predictions at
larger scales~\cite{schwadeEtAl2026}. The targets and representations are
method-specific, but both examples reinforce why a matrix-level objective must
name its representation rather than inherit meaning from its eigenvalues.

Some fitting studies augment energy residuals with band gaps, effective masses,
or other observables. Such an objective must state each observable's definition,
weight, and scale separately. A gap is derived from selected eigenvalues, an
effective mass requires a derivative or local fit, and an optical transition
strength additionally depends on states and another operator. These terms
should not be called ``purely spectral'' merely because the final reported
quantity is a scalar. Klimeck and collaborators provide a concrete mixed-target
example: they fitted silicon nearest- and second-neighbor $sp^3s^{\ast}$ models
to band-edge and effective-mass targets, reported a successful second-neighbor
parameterization, and reported that their search found no satisfactory global
nearest-neighbor fit~\cite{klimeck2000}. That last statement is evidence about
the completed search, not by itself an infeasibility proof.

Two matrices can have identical eigenvalues and a nonzero operator residual.
Equation~\eqref{eq:admissible-spectral-rms} therefore cannot replace
Equation~\eqref{eq:admissible-operator-rms}. The two losses ask different
questions.

\section{Alignment variables are not model parameters}
\index{alignment!constrained family}
\index{nuisance parameter}

A composite retained space permits unitary changes of frame. The alignment
matrix $C$ is consequently an additional comparison variable, not necessarily
a physical model parameter. Let $\mathcal U$ be the declared family of allowed
alignments. Assume that $\mathcal U$ is finite, or that it is compact and the
loss depends continuously on $C$, so that the best permitted value is attained.
Define
\begin{equation}
  L_H^{\min}(\boldsymbol\theta)
  =
  \min_{C\in\mathcal U}L_H(\boldsymbol\theta,C).
  \label{eq:admissible-minimized-operator-loss}
\end{equation}
Without an attainment condition, ``minimum'' must be replaced by ``infimum,''
and a boundary value equal to that infimum need not supply an actual alignment
witness. The family $\mathcal U$ is part of the contract. It may be one fixed identity,
a finite list of global rotations, a continuous symmetry family, or some other
physically justified set. These choices are not interchangeable.

Chapter~\ref{ch:training-composite-band-gauge-alignment} distinguished
pointwise Procrustes alignment from one globally constrained unitary. A
pointwise map $C(k_j)$ may recover a known gauge attack exactly at every sample,
while one constant $C$ may not. Failure within one finite family of constant
rotations does not prove failure for every nonidentity alignment, every smooth
gauge, or every physically admissible map. Enlarging $\mathcal U$ changes the
model-comparison problem.

If $\mathcal U=\{C_{\alpha}:\alpha\in\mathcal F\}$ is a finite family, then
Equation~\eqref{eq:admissible-minimized-operator-loss} is evaluated by explicit
enumeration. Such enumeration is deterministic and complete for
$\mathcal F$; it is not a continuous alignment optimizer.

\section{Admissible sets and their intersection}
\index{admissible set!spectral}
\index{admissible set!operator}

Define the parameter-space spectral-admissible set by
\begin{equation}
  \mathcal A_E(\tau_E)
  =
  \left\{
    \boldsymbol\theta\in\Theta:
    L_E(\boldsymbol\theta)\leq\tau_E
  \right\}
  \label{eq:admissible-spectral-set}
\end{equation}
and the operator-admissible set by
\begin{equation}
  \mathcal A_H(\tau_H)
  =
  \left\{
    \boldsymbol\theta\in\Theta:
    L_H^{\min}(\boldsymbol\theta)\leq\tau_H
  \right\}.
  \label{eq:admissible-operator-set}
\end{equation}
For a finite alignment family,
\begin{equation}
  \mathcal A_H(\tau_H)
  =
  \bigcup_{\alpha\in\mathcal F}
  \left\{
    \boldsymbol\theta\in\Theta:
    L_H(\boldsymbol\theta,C_{\alpha})\leq\tau_H
  \right\}.
  \label{eq:admissible-operator-union}
\end{equation}
Thus the operator-admissible set need not be one connected region even when
each fixed-alignment loss has a simple shape.

The family is compatible at the declared thresholds when
\begin{equation}
  \mathcal A_E(\tau_E)
  \cap
  \mathcal A_H(\tau_H)
  \neq\varnothing.
  \label{eq:admissible-common-intersection}
\end{equation}
Any retained $\boldsymbol\theta_{\mathrm w}$ satisfying both inequalities is a
common witness. A complete operator witness also records an alignment
$C_{\mathrm w}\in\mathcal U$ that attains the accepted operator loss.
Substitution into independently evaluated losses proves membership directly.

A witness proves existence but does not establish uniqueness, parameter
identifiability, robustness, or physical validity. It also says nothing about
parameters outside the frozen family.

This intersection is the contribution boundary. Spectral fitting, direct
Hamiltonian regression, and analytic restrictions on particular tight-binding
classes are established ideas. The question introduced here is narrower:
whether one candidate in a frozen family satisfies independently declared
spectral and aligned-operator criteria. The answer must be supported by a
common witness or, when the sets do not intersect, by a valid positive
separation certificate under the declared family, alignment, losses, metric,
and thresholds.

\section{Why quadratic losses produce ellipses}
\index{quadratic loss}
\index{ellipse!admissible set}

The geometry becomes especially transparent when a residual depends linearly
on two parameters. Suppose
\begin{equation}
  r_{\ell}(\boldsymbol\theta)
  =
  \mathbf a_{\ell}^{\mathsf T}\boldsymbol\theta-b_{\ell},
  \qquad \ell=1,\ldots,M,
\end{equation}
and define
\begin{equation}
  L^2(\boldsymbol\theta)
  =
  \frac{1}{M}
  \sum_{\ell=1}^{M}
  r_{\ell}(\boldsymbol\theta)^2.
\end{equation}
Expanding the squares gives
\begin{equation}
  L^2(\boldsymbol\theta)
  =
  \boldsymbol\theta^{\mathsf T}Q\boldsymbol\theta
  +2\mathbf g^{\mathsf T}\boldsymbol\theta+d,
  \label{eq:admissible-expanded-quadratic}
\end{equation}
where
\begin{equation}
  Q=\frac{1}{M}\sum_{\ell=1}^{M}
      \mathbf a_{\ell}\mathbf a_{\ell}^{\mathsf T},
  \qquad
  \mathbf g=-\frac{1}{M}\sum_{\ell=1}^{M}b_{\ell}\mathbf a_{\ell},
  \qquad
  d=\frac{1}{M}\sum_{\ell=1}^{M}b_{\ell}^{2}.
\end{equation}
If $Q$ is positive definite, the center is
\begin{equation}
  \mathbf c=-Q^{-1}\mathbf g.
\end{equation}
Completing the square yields
\begin{equation}
  L^2(\boldsymbol\theta)
  =
  L_{\min}^{2}
  +
  (\boldsymbol\theta-\mathbf c)^{\mathsf T}
  Q
  (\boldsymbol\theta-\mathbf c),
  \label{eq:admissible-centered-quadratic}
\end{equation}
with
\begin{equation}
  L_{\min}^{2}=d-\mathbf g^{\mathsf T}Q^{-1}\mathbf g.
\end{equation}
Consequently, the unconstrained threshold set $L\leq\tau$ is empty if
$\tau^2<L_{\min}^2$, is the single point $\mathbf c$ if
$\tau^2=L_{\min}^2$, and is a filled ellipse if
$\tau^2>L_{\min}^2$. Intersecting that ellipse with a rectangular parameter
domain may clip its boundary. If $Q$ is only positive semidefinite, flat
directions require separate treatment; the positive-definite formulas must not
be used without checking their premise.

For a finite alignment family, each fixed-alignment operator loss may generate
one ellipse. Equation~\eqref{eq:admissible-operator-union} then becomes a union
of clipped ellipses.

\subsection{Example 2: a witness and a separation distance in two dimensions}

Let $\boldsymbol\theta=(s,q)^{\mathsf T}$ and consider the circular quadratic
losses
\begin{align}
  L_E^2(s,q)&=s^2+(q-0.9)^2,\\
  L_H^2(s,q)&=(s-0.2)^2+(q-0.5)^2.
\end{align}
Their centers are
\begin{equation}
  \mathbf c_E=
  \begin{bmatrix}0\\0.9\end{bmatrix},
  \qquad
  \mathbf c_H=
  \begin{bmatrix}0.2\\0.5\end{bmatrix},
\end{equation}
and the center distance is
\begin{equation}
  d_c
  =
  \lVert\mathbf c_E-\mathbf c_H\rVert_2
  =
  \sqrt{0.2}
  \approx 0.447214.
\end{equation}
With equal thresholds $0.3$, the midpoint
\begin{equation}
  \boldsymbol\theta_{\mathrm w}
  =
  \frac{\mathbf c_E+\mathbf c_H}{2}
  =
  \begin{bmatrix}0.1\\0.7\end{bmatrix}
\end{equation}
has both losses equal to $d_c/2\approx0.223607$. It is therefore a common
witness.

With equal thresholds $0.2$, the two closed disks are disjoint. Their exact
Euclidean separation is
\begin{equation}
  \delta^{\ast}=d_c-0.2-0.2\approx0.047214.
\end{equation}
This exact expression is available because the sets are ordinary disks. General
ellipses and unions require more care.

\begin{pythonexample}{checking the two-dimensional witness and separation}
import numpy as np

# These centers and radii come from the symbolic calculation above.
center_e = np.array([0.0, 0.9])
center_h = np.array([0.2, 0.5])
center_distance = np.linalg.norm(center_e - center_h)

# At radius 0.3, the midpoint is an explicit common witness.
witness = 0.5 * (center_e + center_h)
loss_e = np.linalg.norm(witness - center_e)
loss_h = np.linalg.norm(witness - center_h)

# At radius 0.2, the disk-to-disk distance is positive.
separation = max(0.0, center_distance - 0.2 - 0.2)

print("center distance:", center_distance)
print("candidate witness:", witness)
print("spectral loss at witness:", loss_e)
print("operator loss at witness:", loss_h)
print("witness passes radius 0.3:", loss_e <= 0.3 and loss_h <= 0.3)
print("separation at radius 0.2:", separation)
\end{pythonexample}

\section{Witnesses and separation certificates}
\index{witness!common admissible}
\index{separation certificate}

A common witness is easy to verify: evaluate both losses at the same retained
parameter and compare them with the frozen thresholds. Proving that no witness
exists is harder.

Let $d_{\Theta}$ be a declared metric on the parameter domain. Define
\begin{equation}
  \delta^{\ast}
  =
  \inf_{\substack{
    \boldsymbol\theta_E\in\mathcal A_E\\
    \boldsymbol\theta_H\in\mathcal A_H}}
  d_{\Theta}(\boldsymbol\theta_E,\boldsymbol\theta_H).
  \label{eq:admissible-set-separation}
\end{equation}
If either admissible set is empty, that fact should be reported separately
rather than hidden inside an arbitrary distance convention. If both are
nonempty, a certified lower bound
\begin{equation}
  0<\delta_{\mathrm{lower}}\leq\delta^{\ast}
\end{equation}
proves separation. A displayed pair of feasible boundary points supplies an
upper bound
\begin{equation}
  \delta^{\ast}
  \leq
  d_{\Theta}(\boldsymbol\theta_E^{\mathrm b},
             \boldsymbol\theta_H^{\mathrm b})
  =\delta_{\mathrm{upper}}.
\end{equation}
Equal lower and upper bounds identify the distance exactly, up to the declared
numerical tolerance. If the protocol declares a separation resolution
$\rho>0$, the disposition is certified separated only when the valid lower
bound exceeds $\rho$. A positive bound no larger than $\rho$ remains unresolved
at that declared resolution.

For Euclidean coordinates, one useful certificate projects both sets onto a
unit direction $\mathbf e$. If every spectral-admissible point satisfies
\begin{equation}
  \mathbf e^{\mathsf T}\boldsymbol\theta_E\geq a
\end{equation}
while every operator-admissible point satisfies
\begin{equation}
  \mathbf e^{\mathsf T}\boldsymbol\theta_H\leq b,
\end{equation}
then
\begin{equation}
  \delta^{\ast}\geq\max(0,a-b).
  \label{eq:admissible-directional-lower-bound}
\end{equation}
Both inequalities must hold over the complete applicable sets, including every
alignment branch and every effect of clipping by $\Theta$.

For the ellipse in Equation~\eqref{eq:admissible-centered-quadratic}, its
unclipped extreme in direction $\mathbf e$ is
\begin{equation}
  \mathbf e^{\mathsf T}\mathbf c
  \mathbin{\pm}
  \sqrt{
    (\tau^2-L_{\min}^2)
    \mathbf e^{\mathsf T}Q^{-1}\mathbf e
  }.
  \label{eq:admissible-ellipse-directional-extreme}
\end{equation}
This formula follows by transforming the ellipse to a unit ball and applying
the Cauchy--Schwarz inequality. It is valid only when the ellipse is nonempty,
$Q$ is positive definite, and the relevant extreme is not changed by the
parameter-domain boundary. Those conditions are premises of the certificate,
not implementation details.

A local optimizer, a dense plot, or a finite grid may suggest separation but
does not establish the global lower bound in
Equation~\eqref{eq:admissible-set-separation}. The nearest-neighbor search
reported by Klimeck and collaborators above is an instructive example of this
evidence boundary: ``no satisfactory fit found'' does not become a theorem
that the admissible set is empty. Validated global-search methods can sometimes
provide such bounds, but their applicability and assumptions must be stated
explicitly~\cite{neumaier2004}.

\subsection{The parameter metric is part of the claim}

Euclidean distance is not invariant under arbitrary reparameterization. If one
coordinate is replaced by $100$ times itself, the numerical distance changes
unless the metric changes too. A separation result must therefore state
\begin{itemize}
  \item the canonical parameter coordinates;
  \item their units or dimensionless normalization;
  \item the metric and any coordinate weights; and
  \item the resolution below which separation is treated as unresolved.
\end{itemize}
A positive distance in a synthetic parameter metric is not automatically a
physical energy, probability, or uncertainty interval.

\section{Training, withheld evaluation, and locality}
\index{training data!admissible set}
\index{withheld validation!admissible set}
\index{locality!diagnostic role}

Training data define the losses whose sublevel sets enter the compatibility
decision. Withheld data evaluate selected parameters after that decision. If a
withheld mesh is consulted while choosing thresholds, changing the family, or
selecting the witness, then it is no longer purely withheld for that decision.

For selected $\boldsymbol\theta$, one may report
\begin{align}
  L_E^{\mathrm{train}}(\boldsymbol\theta),
  &\qquad
  L_H^{\mathrm{train}}(\boldsymbol\theta),\\
  L_E^{\mathrm{withheld}}(\boldsymbol\theta),
  &\qquad
  L_H^{\mathrm{withheld}}(\boldsymbol\theta).
\end{align}
The withheld values test interpolation or generalization to the frozen
withheld coordinates. They do not retroactively alter the training-set
admissible sets unless the protocol explicitly defines a second acceptance
stage.

A reciprocal operator may also be transformed to hopping blocks and truncated
at several real-space ranges. The omitted-block norm then measures locality in
the chosen gauge and convention. Unless a locality threshold appears in the
formal definition of admissibility, this quantity is a diagnostic rather than
a third acceptance condition. Reporting it beside the two losses must not
silently turn it into one.

\section{Threshold sensitivity}
\index{threshold sensitivity}

Compatibility is relative to $(\tau_E,\tau_H)$. Increasing a threshold can
enlarge its admissible set; decreasing it can shrink the set or make it empty.
A threshold-sensitivity calculation asks how the disposition changes as one or
both thresholds vary.

Such an analysis is useful for identifying
\begin{itemize}
  \item whether a conclusion is far from or close to a transition;
  \item which criterion controls the transition;
  \item whether a reported positive separation is numerically resolved; and
  \item whether small changes in a proposed threshold change the disposition.
\end{itemize}
It does not calibrate a physically meaningful threshold by itself. If the
sensitivity sweep was designed after the confirmatory result was inspected, it
must be labelled post hoc. It may characterize the retained geometry without
becoming a prospectively frozen prediction.

\section{The controlled M3 role}
\index{M3 controlled calculation}

The maintained one-dimensional M3 calculation instantiates the framework above
with a synthetic rank-two parent inherited from the controlled multiband
calculation, the two parameters $(s,q)$ from
Equation~\eqref{eq:admissible-rank-two-family}, a bounded rectangular domain,
and a frozen finite family of global real rotations. It uses uniform sample
weights as the special case $w_j=1$, constructs spectral and operator quadratic
losses on the training mesh, evaluates selected points
on a disjoint withheld mesh, and records hopping-locality diagnostics at
selected ranges.

This role is deliberately narrow. M3 is a controlled example of how to retain
one common witness or a positive analytic separation certificate within one
frozen family. It is not a general continuous alignment solver, a statement
about every rank-two Hamiltonian, a localization calculation, a material
calculation, or a calibration of physical acceptance tolerances. The retained
outcomes and their provenance belong in
Section~\ref{sec:appendix-g-m3} and the owning calculation package, not in this
foundational derivation. Figure~\ref{fig:periodic-1d-m3-admissible-sets} there
shows the actual compatible and separated M3 regions; the circles in
Example~2 above remain a simpler pedagogical model.

\section{Common mistakes}

\begin{enumerate}
  \item \textbf{Presenting the loss as the contribution.} Weighted spectral
        fitting, effective-mass targets, and direct Hamiltonian regression all
        have precedents. The present contribution is the bounded simultaneous
        witness-or-certificate question.
  \item \textbf{Comparing two separate minimizers.} The spectral optimum and
        operator optimum need not define one common candidate.
  \item \textbf{Calling a failed grid search a proof.} A grid leaves untested
        points between samples and supplies no global lower bound by itself.
  \item \textbf{Optimizing over an undeclared gauge family.} Changing the
        alignment family changes the admissible set.
  \item \textbf{Using incompatible represented operators.} A Frobenius norm
        cannot repair mismatched spaces, bases, units, or energy references.
  \item \textbf{Choosing thresholds from the desired answer.} Thresholds are
        acceptance inputs, not knobs adjusted until an intersection appears or
        disappears.
  \item \textbf{Treating withheld values as training evidence.} Withheld data
        have a distinct evaluation role.
  \item \textbf{Promoting a locality diagnostic to a criterion.} That requires
        a declared locality loss and threshold.
  \item \textbf{Generalizing beyond the frozen family.} Separation inside one
        bounded family does not establish universal incompatibility.
\end{enumerate}

\section{Exercises}

\begin{enumerate}
  \item For Example~1, determine the intersection for equal thresholds $0.5$.
        Is there one witness or more than one?
  \item Starting from Equation~\eqref{eq:admissible-expanded-quadratic}, carry
        out the completion of the square and verify
        Equation~\eqref{eq:admissible-centered-quadratic}.
  \item Let $Q=\operatorname{diag}(4,1)$, $\mathbf c=\mathbf0$, and
        $L_{\min}=0$. Sketch the sets for $\tau=1$ and $\tau=1/2$ and explain
        which parameter is more tightly constrained.
  \item Add a second operator-alignment branch centered at $(-0.3,0.4)$ to
        Example~2. Draw the union of the two operator disks and determine which
        branch is closer to the spectral center.
  \item Explain why multiplying the first parameter coordinate by $100$
        changes an unweighted Euclidean separation. Propose a dimensionless
        metric that preserves the intended relative scales.
  \item Construct an example in which no sampled grid point satisfies both
        criteria but an unsampled common witness exists.
  \item State which additional contract would be required to make an
        omitted-hopping norm part of the admissibility decision rather than a
        diagnostic.
\end{enumerate}

\section{Chapter summary}

Spectral fitting, direct Hamiltonian regression, and analytic restrictions on
particular tight-binding classes predate this framework. The role of the
framework is to make a narrower simultaneous-adequacy question explicit.
A model family is compatible with two declared criteria only when the same
parameter lies in both thresholded admissible sets. A common witness proves
that such a parameter exists. A positive, globally valid lower bound proves
separation; failure of a finite search does not. Quadratic losses make this
logic visible through ellipses and finite unions of ellipses, but their analytic
certificates remain conditional on positive curvature, domain geometry, and
the declared alignment family. Training data define the decision, withheld
data evaluate selected candidates, and locality remains a diagnostic unless it
is explicitly promoted to a thresholded criterion. These distinctions allow
later bulk and impurity chapters to use the language of adequacy without
turning a bounded synthetic calculation into a material claim.
